\chapter{Conclusion générale et perspectives}

\section{Bilan du projet}

Le projet SIRH-Doc a permis de concevoir une plateforme complète autour d'un besoin administratif très concret : produire des documents RH fiables à partir de données centralisées et traçables. Au départ, le sujet pouvait sembler proche d'un simple éditeur de modèles. En pratique, il a rapidement demandé de traiter des questions plus profondes : qui a le droit de voir quelles données, dans quelle organisation, pour quelle année universitaire, et avec quelle preuve après génération.

L'un des apports majeurs du projet réside dans l'approche metadata-driven. Les modules, les vues de tables, les filtres, les familles de documents et les variables ne sont pas figés dans le code. Ils sont configurés, puis interprétés par l'application. Ce choix a parfois compliqué le développement, notamment pour les requêtes SQL dynamiques et les variables de type tableau, mais il correspond mieux à un environnement où les besoins changent d'une organisation à l'autre.

\section{Apports techniques}

Sur le plan backend, le projet a consolidé une architecture ASP.NET Core organisée autour de contrôleurs, services et repositories. Dapper a permis de conserver une maîtrise directe des requêtes SQL Server, ce qui s'est révélé utile pour les lectures dynamiques, les filtres d'année universitaire et les requêtes de génération documentaire. Le middleware de résolution tenant constitue une pièce centrale : il assure que chaque opération s'effectue dans le bon contexte.

Côté frontend, Angular a fourni une base solide pour construire des interfaces administratives denses : super administration, modules, vues de tables, studio documentaire, génération et archives. Les services Angular isolent les appels API et la normalisation, tandis que les composants se concentrent sur les interactions utilisateur.

\section{Limites rencontrées}

Certaines limites restent présentes. La pagination documentaire en HTML reste un sujet délicat, car les navigateurs ne sont pas des moteurs de composition dédiés à l'édition papier. La gestion des images, des tableaux longs et des ruptures de page demande encore des tests sur un plus grand nombre de documents réels. C'est typiquement le genre de point qui ne se règle pas uniquement par une bonne architecture : il faut aussi observer les cas concrets, ajuster les styles et accepter plusieurs itérations.

Le moteur dynamique, bien qu'efficace, exige une configuration rigoureuse. Une mauvaise définition de filtre, de lookup ou de requête SQL peut produire un comportement inattendu. Il serait donc utile d'ajouter davantage d'outils de diagnostic pour les super administrateurs : prévisualisation SQL, validation automatique des colonnes et rapports de cohérence.

\section{Perspectives}

Plusieurs évolutions peuvent être envisagées. La première concerne la signature électronique, qui donnerait une valeur juridique plus forte aux documents générés. La deuxième porte sur l'envoi automatique par courriel, notamment pour les attestations ou documents destinés à des agents externes. Une troisième perspective serait l'intégration d'un workflow de validation : un document pourrait être préparé par un utilisateur, vérifié par un responsable, puis archivé après validation.

À moyen terme, le système pourrait également intégrer des tableaux de bord permettant de suivre le volume de documents générés, les familles les plus utilisées, les délais de traitement ou les erreurs fréquentes. Ces indicateurs aideraient l'administration à piloter plus finement son activité documentaire.

\section{Conclusion}

Au final, ce projet a permis de traiter un problème d'ingénierie complet : comprendre un besoin métier, concevoir une architecture sécurisée, construire un moteur configurable, produire une interface exploitable et assurer la traçabilité des opérations. SIRH-Doc constitue aujourd'hui une base opérationnelle crédible pour accompagner la dématérialisation RH d'un établissement éducatif. Il reste perfectible, mais ses fondations permettent déjà d'envisager des extensions sans remettre en cause l'ensemble du système.
